% Figure Legends (250 words per legend)
% - Should always have a short, explanatory title (as well as legend).
% - Legends must explain the figure fully without reference to the text.
% - If applicable, the original source of previously published material must be acknowledged in the Figure legend.
% - Figures must be cited in the main text (box figures must be cited in the box). 

\begin{figure}[p]
  \centering
\includegraphics[width=14cm]{figures/environments}
\caption{\textbf{Examples of games and environments for emergent communication.}
(a) Emergent communication work in the pre-deep-learning era typically used symbolic data as input: \citeA{Batali:1998} presents a study where recurrent neural network agents communicate in a referential game using sequences of discrete symbols. Similar work with deep networks often uses realistic pictures as input, see Fig.~\ref{fig:referential-game} for an example. (b) More complex scenarios with deep agents: \citeA{Cao:etal:2018} study self-interested agents engaging in a multi-turn negotiation game. (c) Richer, dynamic environments: \citeA{Jaques:etal:2019} study five embodied self-interested agents engaging in multi-turn interactions while navigating in a 2D visual environment. (d) Scaling up to fully realistic scenarios: in the experiment of \citeA{Das:etal:2019}, embodied cooperative agents solve navigation challenges in a 3D environment. Images from \citeA{Jaques:etal:2019} and \citeA{Das:etal:2019} reproduced by permission.}\label{fig:environments}
\end{figure}
